\documentclass{article}
\usepackage{amsmath,amssymb,amsthm}
\usepackage{algorithm}
\usepackage{algorithmic}
\usepackage{enumitem}
\usepackage{hyperref}
\newtheorem{theorem}{Theorem}
\newtheorem{lemma}{Lemma}
\newcommand{\prox}{\operatorname{prox}}
\begin{document}

\section*{Proximal Gradient Method (PGM)}

\section*{Mathematical Formulation}
Let $F(x)=f(x)+g(x)$ where
\begin{itemize}
  \item $f:\mathbb{R}^{d}\rightarrow\mathbb{R}$ is convex and differentiable with $L$‑Lipschitz continuous gradient, i.e. 
        $\|\nabla f(x)-\nabla f(y)\|\le L\|x-y\|,\;\forall x,y$.
  \item $g:\mathbb{R}^{d}\rightarrow\mathbb{R}\cup\{+\infty\}$ is proper, closed and convex (possibly nonsmooth).
\end{itemize}
Given a stepsize $\eta>0$, the proximal gradient iteration is
\begin{equation}
\label{eq:pgm}
x_{k+1}= \prox_{\eta g}\!\bigl(x_k-\eta \nabla f(x_k)\bigr),\qquad k=0,1,2,\dots
\end{equation}
where the proximal operator of $g$ is
\[
\prox_{\eta g}(v)=\arg\min_{u\in\mathbb{R}^{d}}\Bigl\{ g(u)+\frac{1}{2\eta}\|u-v\|^{2}\Bigr\}.
\]

\section*{Pseudocode}
\begin{algorithm}[H]
\caption{Proximal Gradient Method}
\begin{algorithmic}[1]
\REQUIRE Initial point $x_{0}\in\mathbb{R}^{d}$, stepsize $\eta\in(0,1/L]$.
\FOR{$k=0,1,2,\dots$}
    \STATE $v_{k}=x_{k}-\eta\nabla f(x_{k})$ \COMMENT{gradient step}
    \STATE $x_{k+1}= \prox_{\eta g}(v_{k})$ \COMMENT{proximal step}
    \IF{termination criterion satisfied}
        \STATE \textbf{break}
    \ENDIF
\ENDFOR
\RETURN $x_{k+1}$
\end{algorithmic}
\end{algorithm}

\section*{Dry Run Example}
Consider the scalar problem
\[
f(w)=(w-3)^{2},\qquad g(w)\equiv0,\qquad \eta=0.1,\qquad w_{0}=0.
\]
Since $g\equiv0$, $\prox_{\eta g}(v)=v$.
\begin{enumerate}
  \item $k=0$: $\nabla f(w_{0})=2(w_{0}-3)=-6$.
        \[
        v_{0}=w_{0}-\eta\nabla f(w_{0})=0-0.1(-6)=0.6,\qquad
        w_{1}= \prox_{\eta g}(v_{0})=0.6.
        \]
        Objective $F(w_{1})=f(0.6)=(0.6-3)^{2}=5.76$.
  \item $k=1$: $\nabla f(w_{1})=2(0.6-3)=-4.8$.
        \[
        v_{1}=0.6-0.1(-4.8)=1.08,\qquad
        w_{2}=1.08,
        \]
        $F(w_{2})=f(1.08)=(1.08-3)^{2}=3.6864$.
  \item $k=2$: $\nabla f(w_{2})=2(1.08-3)=-3.84$.
        \[
        v_{2}=1.08-0.1(-3.84)=1.464,\qquad
        w_{3}=1.464,
        \]
        $F(w_{3})=f(1.464)=(1.464-3)^{2}=2.3660$.
\end{enumerate}
The iterates move toward the minimizer $w^{*}=3$ and the objective values decrease.

\newpage

\section*{Assumptions}
\begin{enumerate}[label={A\arabic*}.]
  \item \textbf{Convexity of $f$ and $g$}: $f$ and $g$ are convex; $f$ is differentiable.
  \item \textbf{Lipschitz Gradient}: $\nabla f$ is $L$‑Lipschitz:
        $\|\nabla f(x)-\nabla f(y)\|\le L\|x-y\|,\;\forall x,y$.
  \item \textbf{Stepsize}: $\eta\in(0,1/L]$.
  \item \textbf{Existence of a Minimizer}: The set 
        $\mathcal{X}^{*}=\arg\min_{x}F(x)$ is non‑empty; denote any $x^{*}\in\mathcal{X}^{*}$.
\end{enumerate}

\section*{Main Convergence Theorem}
\begin{theorem}
\label{thm:rate}
Under Assumptions A1–A4 and the iteration \eqref{eq:pgm} with $\eta\in(0,1/L]$, the sequence $\{x_{k}\}$ satisfies for all $k\ge1$,
\begin{equation}
\label{eq:rate}
F(x_{k})-F(x^{*})\;\le\;\frac{\|x_{0}-x^{*}\|^{2}}{2\eta k}.
\end{equation}
Consequently $F(x_{k})\to F(x^{*})$ at the rate $O(1/k)$.
\end{theorem}

\section*{Theoretical Properties}
\begin{itemize}
  \item \textbf{Stability}: The iteration is a non‑expansive mapping because the proximal operator is firmly non‑expansive and the gradient step is $\frac{1}{L}$‑Lipschitz; thus the whole map is $\frac{1}{1-\eta L}$‑Lipschitz $\le1$.
  \item \textbf{Hyperparameter Sensitivity}: The stepsize must respect $\eta\le1/L$; larger $\eta$ may violate the descent inequality (see Lemma~\ref{lem:descent}).
  \item \textbf{Comparison to Gradient Descent}: When $g\equiv0$, \eqref{eq:rate} reduces to the classical $O(1/k)$ rate for gradient descent on smooth convex $f$.
  \item \textbf{Extension}: With strong convexity of $F$, a linear rate $O((1-\mu\eta)^{k})$ can be derived (not shown).
\end{itemize}

\section*{Discussion}
The theorem shows that even with a nonsmooth convex regularizer $g$, the simple proximal gradient scheme retains the optimal $O(1/k)$ convergence rate for composite convex optimization. The proof hinges on two elementary tools: (i) the descent lemma for $L$‑smooth $f$, and (ii) the non‑expansiveness (firm non‑expansiveness) of the proximal operator. No stochasticity or variance reduction is required.

\newpage

\section*{Preliminaries (Definitions and Lemmas)}
\begin{definition}[Proximal Operator]
For $\eta>0$, $\prox_{\eta g}(v)=\arg\min_{u}\bigl\{g(u)+\frac{1}{2\eta}\|u-v\|^{2}\bigr\}$.
\end{definition}
\begin{lemma}[Firm Non‑expansiveness of $\prox$]
\label{lem:firm}
For any $u,v\in\mathbb{R}^{d}$,
\[
\|\prox_{\eta g}(u)-\prox_{\eta g}(v)\|^{2}
\le\langle \prox_{\eta g}(u)-\prox_{\eta g}(v),\,u-v\rangle .
\]
Consequently,
\[
\|\prox_{\eta g}(u)-\prox_{\eta g}(v)\|\le\|u-v\|.
\]
\end{lemma}
\begin{lemma}[Descent Lemma for $L$‑smooth $f$]
\label{lem:descent}
For any $x,y\in\mathbb{R}^{d}$,
\[
f(y)\le f(x)+\langle\nabla f(x),y-x\rangle+\frac{L}{2}\|y-x\|^{2}.
\]
\end{lemma}

\section*{Main Proof (Step‑by‑Step Derivation)}
Let $x^{*}\in\mathcal{X}^{*}$ be any optimal solution. Define the gradient step
\[
v_{k}=x_{k}-\eta\nabla f(x_{k}) .
\]
The proximal update is $x_{k+1}= \prox_{\eta g}(v_{k})$.

\begin{enumerate}
\item[Step 1.] \textbf{Optimality condition of the proximal step.}  
By definition of $\prox$, $x_{k+1}$ satisfies
\[
0 \in \partial g(x_{k+1})+\frac{1}{\eta}(x_{k+1}-v_{k}),
\]
i.e.
\begin{equation}
\label{eq:optcond}
\frac{1}{\eta}(v_{k}-x_{k+1})\in\partial g(x_{k+1}).
\end{equation}
\item[Step 2.] \textbf{Convex subgradient inequality for $g$.}  
Using \eqref{eq:optcond} and convexity of $g$,
\[
g(x^{*})\ge g(x_{k+1})+\left\langle \frac{1}{\eta}(v_{k}-x_{k+1}),\,x^{*}-x_{k+1}\right\rangle .
\tag{2.1}
\]
\item[Step 3.] \textbf{Apply the descent lemma to $f$.}  
From Lemma~\ref{lem:descent} with $x=x_{k}$ and $y=x_{k+1}$,
\[
f(x_{k+1})\le f(x_{k})+\langle\nabla f(x_{k}),x_{k+1}-x_{k}\rangle+\frac{L}{2}\|x_{k+1}-x_{k}\|^{2}.
\tag{3.1}
\]
\item[Step 4.] \textbf{Combine $f$ and $g$ inequalities.}  
Add (2.1) and (3.1):
\[
\begin{aligned}
F(x_{k+1})-F(x^{*})
&\le \langle\nabla f(x_{k}),x_{k+1}-x_{k}\rangle
   +\frac{L}{2}\|x_{k+1}-x_{k}\|^{2}
   -\frac{1}{\eta}\langle v_{k}-x_{k+1},\,x^{*}-x_{k+1}\rangle .
\end{aligned}
\tag{4.1}
\]
\item[Step 5.] \textbf{Express $v_{k}$ in terms of $x_{k}$.}  
Recall $v_{k}=x_{k}-\eta\nabla f(x_{k})$. Hence
\[
v_{k}-x_{k+1}=x_{k}-\eta\nabla f(x_{k})-x_{k+1}.
\]
Insert into (4.1):
\[
\begin{aligned}
F(x_{k+1})-F(x^{*})
&\le \langle\nabla f(x_{k}),x_{k+1}-x_{k}\rangle
   +\frac{L}{2}\|x_{k+1}-x_{k}\|^{2}\\
&\qquad -\frac{1}{\eta}\langle x_{k}-\eta\nabla f(x_{k})-x_{k+1},\,x^{*}-x_{k+1}\rangle .
\end{aligned}
\tag{5.1}
\]
\item[Step 6.] \textbf{Expand the inner product.}  
\[
\begin{aligned}
-\frac{1}{\eta}\langle x_{k}-x_{k+1},x^{*}-x_{k+1}\rangle
+\langle\nabla f(x_{k}),x^{*}-x_{k+1}\rangle .
\end{aligned}
\tag{6.1}
\]
Thus (5.1) becomes
\[
\begin{aligned}
F(x_{k+1})-F(x^{*})
&\le \frac{L}{2}\|x_{k+1}-x_{k}\|^{2}
   -\frac{1}{\eta}\langle x_{k}-x_{k+1},x^{*}-x_{k+1}\rangle .
\end{aligned}
\tag{6.2}
\]
(The terms $\langle\nabla f(x_{k}),x_{k+1}-x_{k}\rangle$ cancel with the last term of (6.1).)
\item[Step 7.] \textbf{Use the identity $2\langle a,b\rangle =\|a\|^{2}+\|b\|^{2}-\|a-b\|^{2}$.}  
Set $a=x_{k}-x_{k+1}$ and $b=x_{k+1}-x^{*}$:
\[
\langle x_{k}-x_{k+1},x^{*}-x_{k+1}\rangle
= -\frac12\Bigl(\|x_{k}-x^{*}\|^{2}
               -\|x_{k+1}-x^{*}\|^{2}
               -\|x_{k}-x_{k+1}\|^{2}\Bigr).
\tag{7.1}
\]
Plug (7.1) into (6.2):
\[
\begin{aligned}
F(x_{k+1})-F(x^{*})
&\le \frac{L}{2}\|x_{k+1}-x_{k}\|^{2}
   +\frac{1}{2\eta}\Bigl(\|x_{k}-x^{*}\|^{2}
               -\|x_{k+1}-x^{*}\|^{2}
               -\|x_{k}-x_{k+1}\|^{2}\Bigr)\\
&= \frac{1}{2\eta}\bigl(\|x_{k}-x^{*}\|^{2}-\|x_{k+1}-x^{*}\|^{2}\bigr)
   +\Bigl(\frac{L}{2}-\frac{1}{2\eta}\Bigr)\|x_{k+1}-x_{k}\|^{2}.
\end{aligned}
\tag{7.2}
\]
\item[Step 8.] \textbf{Choose stepsize $\eta\le 1/L$.}  
Because $\eta\le 1/L$, we have $L/2-1/(2\eta)\le0$, therefore the second term in (7.2) is non‑positive and can be dropped:
\[
F(x_{k+1})-F(x^{*})\le \frac{1}{2\eta}\bigl(\|x_{k}-x^{*}\|^{2}-\|x_{k+1}-x^{*}\|^{2}\bigr).
\tag{8.1}
\]
\item[Step 9.] \textbf{Telescoping sum.}  
Sum (8.1) for $k=0,\dots,T-1$:
\[
\sum_{k=0}^{T-1}\bigl(F(x_{k+1})-F(x^{*})\bigr)
\le \frac{1}{2\eta}\bigl(\|x_{0}-x^{*}\|^{2}-\|x_{T}-x^{*}\|^{2}\bigr)
\le \frac{\|x_{0}-x^{*}\|^{2}}{2\eta}.
\tag{9.1}
\]
\item[Step 10.] \textbf{Derive the $O(1/k)$ bound.}  
Since each term $F(x_{k+1})-F(x^{*})$ is non‑negative,
\[
T\bigl(F(x_{T})-F(x^{*})\bigr)
\le \sum_{k=0}^{T-1}\bigl(F(x_{k+1})-F(x^{*})\bigr)
\le \frac{\|x_{0}-x^{*}\|^{2}}{2\eta},
\]
which yields
\[
F(x_{T})-F(x^{*})\le \frac{\|x_{0}-x^{*}\|^{2}}{2\eta T}.
\tag{10.1}
\]
\end{enumerate}
Thus Theorem~\ref{thm:rate} follows directly from (10.1).

\section*{Rate Derivation (With Constants)}
The constant in \eqref{eq:rate} is explicitly $\frac{1}{2\eta}$ multiplied by the squared distance of the initial point to the solution set. If we use the maximal admissible stepsize $\eta=1/L$, the bound becomes
\[
F(x_{k})-F(x^{*})\le \frac{L\|x_{0}-x^{*}\|^{2}}{2k}.
\]

\section*{Special Cases}
\begin{itemize}
  \item \textbf{Pure smooth case $g\equiv0$}: $\prox_{\eta g}$ is identity, and the proof reduces to the classic gradient descent analysis.
  \item \textbf{Indicator of a convex set $C$}: $g=\iota_{C}$, $\prox_{\eta g}$ is the Euclidean projection onto $C$. The algorithm becomes projected gradient descent with the same $O(1/k)$ guarantee.
  \item \textbf{Lasso regularizer $g(x)=\lambda\|x\|_{1}$}: $\prox$ is soft‑thresholding, yielding the proximal‑gradient (ISTA) method, again with $O(1/k)$ convergence.
\end{itemize}

\end{document}